% Preamble
\documentclass[10pt, aspectratio=169]{beamer}
	% Packages
	\usepackage[english]{babel}
	\usepackage{lipsum}
	\usepackage{mathptmx}
	\usepackage{xcolor}
	\usepackage{multicol}
	\usepackage{graphics}
	\usepackage{float}
	% Themes
	\usefonttheme[onlymath]{serif}
	\usetheme[nofirafonts]{focus}
	\renewcommand{\thempfootnote}{\arabic{mpfootnote}}
	% Cover
	\title{Microeconometrics Using Stata}
	\subtitle{LINEAR INSTRUMENTAL - VARIABLES REGRESSION}
	\author{
		Jhon R. Ordoñez
		\footnote{\label{1} National University of San Cristóbal de Huamanga}
		\footnote{\label{2} Faculty of Economic, Administrative and Accounting Sciences}
		\footnote{\label{3} Professional School of Economics}
	}
	\date{\today}
% Body
\begin{document}
	% Cover
	\begin{frame}
		\maketitle
	\end{frame}
	% Outline
	\begin{frame}[t]
		\frametitle{Outline}
		%\begin{multicols}{2}
		\tableofcontents
		%\end{multicols}
	\end{frame}
	% Section 1
	\section{Exercise 1}
		% Slide 1
		\begin{frame}
			\frametitle{Exercise 1}
				\begin{block}{Exercise 1}
					Estimate by \textbf{2SLS} the same regression model as in \textcolor{red!50}{\textbf{section 7.4.4}}, with the
					instruments \textcolor{blue}{\textbf{multlc}} and \textcolor{blue}{\textbf{firmsz}}. Compare the 2SLS estimates with OLS
					estimates. Perform a test of endogeneity of \textcolor{blue}{\textbf{hi$\_$empunion}}. Perform a test of
					overidentification. State what you conclude. Throughout this exercise,
					perform inference that is robust to heteroskedasticity.
				\end{block}
		\end{frame}
	% Section 2
	\section{Exercise 2}
		% Slide 1
		\begin{frame}
			\frametitle{Exercise 2}
			\begin{block}{Exercise 2}
				Repeat exercise 1 using optimal \textbf{GMM}.
			\end{block}
		\end{frame}
	% Section 3
	\section{Exercise 3}
		% Slide 1
		\begin{frame}
			\frametitle{Exercise 3}
			\begin{block}{Exercise 3}
				Use the model and instruments of exercise 1. Compare the following
				estimators: \textbf{2SLS}, \textbf{LIML}, and optimal \textbf{GMM }given heteroskedastic errors. For the
				last model, estimate the parameters by using the community-contributed
				\textcolor{blue}{\textbf{ivreg2}} command in addition to \textcolor{blue}{\textbf{ivregress}}.
			\end{block}
		\end{frame}
	% Section 4
	\section{Exercise 4}
		% Slide 1
		\begin{frame}
			\frametitle{Exercise 4}
			\begin{block}{Exercise 4}
				Use the model of exercise 1. Compare \textbf{2SLS} estimates as the instruments
				\textcolor{blue}{\textbf{ssiratio}}, \textcolor{blue}{\textbf{lowincome}}, \textcolor{blue}{\textbf{multlc}}, and \textcolor{blue}{\textbf{firmsz}} are progressively added.
			\end{block}
		\end{frame}
	% Section 5
	\section{Exercise 5}
		% Slide 1
		\begin{frame}
			\frametitle{Exercise 5}
			\begin{block}{Exercise 5}
				Use the model and instruments of exercise 1. Perform appropriate
				diagnostics and tests for weak instruments using the \textbf{2SLS} estimator. State
				what you conclude. Throughout this exercise, perform inference assuming
				errors are i.i.d.
			\end{block}
		\end{frame}
	% Section 6
	\section{Exercise 6}
		% Slide 1
		\begin{frame}
			\frametitle{Exercise 6}
			\begin{block}{Exercise 6}
				Use the model and instruments of exercise 1. Use the community-contributed
				\textcolor{blue}{\textbf{condivreg}} command to perform inference for the \textbf{2SLS} estimator.
				Compare the results with those using conventional asymptotics.
			\end{block}
		\end{frame}
	% Section 7
	\section{Exercise  7}
		% Slide 1
		\begin{frame}
			\frametitle{Exercise 7}
			\begin{block}{Exercise 7}
				Use the model and instruments of exercise 1. Use the community-contributed \textcolor{blue}{\textbf{jive}} command,
				and compare estimates and standard errors from
				the four different variants of \textbf{JIVE} and from optimal \textbf{GMM}. Throughout this
				exercise, perform inference that is robust to heteroskedasticity.
			\end{block}
		\end{frame}
	% Section 8
	\section{Exercise 8}
		% Slide 1
		\begin{frame}
			\frametitle{Exercise 8}
			\begin{block}{Exercise 8}
				Fit the \textbf{3SLS} model of \textcolor{red!50}{\textbf{section 7.10}}, and compare the \textbf{3SLS} coefficient estimates
				and standard errors in the \textcolor{blue}{\textbf{ldrugexp}} equation with those from \textbf{2SLS} estimation
				(with default standard errors).
			\end{block}
		\end{frame}
	% Section 9
	\section{Exercise 9}
		% Slide 1
		\begin{frame}[label = slide 1]
			\frametitle{Exercise 9}
			\begin{block}{Exercise 9}
				This question considers the same earnings-schooling dataset as that
				analyzed in \textcolor{green!75}{\textbf{Cameron and Trivedi}} (\textcolor{green!75}{\textbf{2005}}, 111). The data are in
				\textcolor{blue}{\textbf{mus207klingdata.dta}}. The \textcolor{blue}{\textbf{describe}} command provides descriptions of the
				regressors. There are three endogenous regressors (years of schooling, years
				of work experience, and experience-squared) and three instruments (a
				college proximity indicator, age, and age-squared). Interest lies in the
				coefficient of schooling. Perform appropriate diagnostics and tests for weak
				instruments for the following model. State what you conclude. The
				following commands yield the \textbf{IV} estimator:
				\hyperlink{slide 1 - commands}{\beamergotobutton{see commands}}
			\end{block}
		\end{frame}
	% Section 10
	\section{Exercise 10}
		% Slide 1
		\begin{frame}
			\frametitle{Exercise 10}
			\begin{block}{Exercise 10}
				Use the same dataset as the previous question. Treat only \textcolor{blue}{\textbf{grade76}} as
				endogenous, let \textcolor{blue}{\textbf{exp76}} and \textcolor{blue}{\textbf{expsq76}} be exogenous, and use \textcolor{blue}{\textbf{col4}} as the only
				instrument. Perform appropriate diagnostics and tests for a weak instrument,
				and state what you conclude. Then, use the community-contributed
				\textcolor{blue}{\textbf{condivreg}} command to perform inference, and compare the results with
				those using conventional asymptotics.
			\end{block}
		\end{frame}
	% Section 11
	\section{Exercise 11}
		% Slide 1
		\begin{frame}[t]
			%\frametitle{Exercise 11}
			\begin{block}{Exercise 11}
				When an endogenous variable enters the regression nonlinearly, the obvious
				\textbf{IV} estimator is inconsistent and a modification is needed. Specifically,
				suppose $y_{1} = \beta y_{2}^{2} + u $ and the first-stage equation for
				$y_{2}$ is $y_{2} = \pi_{2}z + v $ where the zero-mean errors and are correlated.
				Here the endogenous regressor appears in the structural equation as $y_{2}^{2}$ rather than $y_{2}$.
				The \textbf{IV} estimator is 
					\begin{equation}
						\hat{\beta}_{IV} = \left(\sum_{i}z_{i}y_{2i}^{2} \right)^{-1}\sum_{i}z_{i}y_{1i}
					\end{equation}
				This can be implemented by a regular \textbf{IV} regression of $y$ on $y_{2}^{2}$ with the instrument
				$z$: regress $y_{2}^{2}$ on $z$ and then regress $y_{1}$ on the first-stage prediction $\widehat{y_{2}^{2}}$. 
				If instead we regress $y_{2}$ on $z$
				at the first stage, giving $\hat{y_{2}}$, and then regress $y_{1}$ on $\left(\hat{y}_{2}\right)^{2}$, an inconsistent
				estimate is obtained. Generate a simulation sample to demonstrate these
				points. Consider whether this example can be generalized to other nonlinear
				models where the nonlinearity is in regressors only, so that 
					\begin{equation}
						y_{1} = \mathbf{g}\left(\mathbf{y}_{2}\right)'\mathbf{\beta} + u
					\end{equation}
				where $\mathbf{g}\left(\mathbf{y}_{2}\right)$ is a nonlinear function of $\mathbf{y}_{2}$.
			\end{block}
		\end{frame}
	% Section 12
	\section{Exercise 12}
		% Slide 1
		\begin{frame}
			\frametitle{Exercise 12}
			\begin{block}{Exercise 12}
				This exercise is based on artificially generated data. Set sample size to
				$1,000$. Set seed to \textcolor{blue}{\textbf{10101}}. Generate a sample using the following \textbf{DGP}: 
				$u\sim N(0,1)$; $e \sim 0.6u + N(0,1)$; $z \sim uniform(0,1)$;
				$y_{2} = z + 0.5z^{2} + 0.0625z^{3} + e$; $x \sim N(0,9)$; $y_{1} = y_{2} + 0.5y_{2}^{2} + x + u$
				Estimate the $y_{1}$ equation by \textbf{OLS}. Theoretically, the estimates of the equation
				are inconsistent. Compare all the slope coefficients with their true values,
				and discuss whether the direction of the bias is as you might expect. \textbf{OLS} is
				inconsistent in this case, yet some coefficients are estimated close to their
				true values; explain why. Run the same regression again with the intercept
				constrained to zero, and comment on any differences that you might observe.
			\end{block}
		\end{frame}
	% Section 13
	\section{Exercise 13}
		% Slide 1
		\begin{frame}
			\frametitle{Exercise 13}
			\begin{block}{Exercise 13}
				Continuing the previous exercise, estimate three variants of the $y_{1}$-equation
				by \textbf{2SLS} using as instruments
					\begin{itemize}
						\item[i)] $x$, $z$, and $z^{2}$ 
						\item[ii)] $x$, $z$, $z^{2}$, and $z^{3}$ 
						\item[iii)] $\hat{y}_{2}$ and $\hat{y}_{2}^{2}$
					\end{itemize}
				where $\hat{y}_{2}$ is the fitted value of $y_{2}$ from the reduced form of $y_{2}$. Are the
				instruments chosen in each case valid and relevant? Which of the three
				estimates would you prefer on theoretical grounds? Explain your answer.
				With reference to the \textbf{2SLS} estimates used earlier, apply the test of
				overidentification where appropriate, and interpret your results.
			\end{block}
		\end{frame}
	% Section 14
	\section{Exercise 14}
		% Slide 1
		\begin{frame}
			\frametitle{Exercise 14}
			\begin{block}{Exercise 14}
				This question continues with the model presented in question $12$ above.
				Modify the specification of the $y_{1}$ equation such that the equation errors are
				heteroskedastic, with skedasticity being a quadratic function of $x$; that is, 
					\begin{equation}
						Var(u) = h(x)\times N(0,1)
					\end{equation}
				Reestimate the $y_{1}$ equation by \textbf{2SLS} and two-step
				\textbf{GMM} using as instruments $x$, $z$, $z^{2}$, $z^{3}$. In both cases, apply the test of
				overidentifying restrictions. Compare the results, and comment on the
				outcome of the test.
			\end{block}
		\end{frame}
	% References
	\begin{frame}[t,allowframebreaks]
		\frametitle{References}
		\bibliography{biblio.bib}
		% Style
		\bibliographystyle{apalike}
		% Authors
		\nocite{cameron2022-I}
		\nocite{cameron2022-II}
	\end{frame}
	% Cover
	\begin{frame}
		\maketitle
	\end{frame}
	% Appendix
	\begin{frame}[label = slide 1 - commands]
		\begin{block}{Commands}
			\begin{figure}[H]
				\centering
				\includegraphics[scale = 0.5]{figures/fig1.PNG}
			\end{figure}
			\hyperlink{slide 1}{\beamergotobutton{back}}
		\end{block}
	\end{frame}
\end{document}